\chapter{Compilation and the Compiler}
\label{ch:compilation}

Everything you have done so far has gone through the evaluator: you type an
expression, Mathilda rewrites it until it stops changing, and prints the result.
That machine is wonderfully general---it will happily add a symbol to a bignum,
thread a function over a ragged list, or apply a rule you wrote a moment ago---but
generality has a price. To evaluate \(x^2 + 1\) at a single machine number, the
evaluator still walks a tree, reads attributes, checks for rules, and dispatches
on the head at every node, allocating and freeing expression objects as it goes.
For one evaluation that cost is invisible. Inside a loop that runs a million
times over plain \texttt{double}s, it is almost all of the running time.

This chapter is about the other gear. \B{Compile} takes a numeric expression and
lowers it, once, to typed \emph{bytecode}\index{compilation!bytecode}---a short
program of monomorphic instructions run by a small register machine that never
allocates an expression, never consults the symbol table, and never asks what
type a value has, because the types were all settled at compile time. The answer
is exactly the evaluator's; the speed is eight to several hundred times
greater\index{Compile!performance}. The same machine runs quietly behind
\B{Table}, \B{Plot}, \B{NIntegrate} and \B{Nest} without your asking, a feature
called \emph{auto-compilation}\index{auto-compilation} that this chapter also
takes apart and shows you how to switch off.

A word to prevent one confusion. This is the \mcode{Compile} \emph{builtin}---a
tool inside the running system---not the business of compiling Mathilda itself
from C source, which is Chapter~\ref{ch:building}. Here nothing touches a C
compiler; the ``compiler'' is a component of the evaluator, and its target is a
bytecode of its own.

\section{From an expression tree to bytecode}
\label{sec:compile-basics}

\B{Compile} takes two arguments: a list of parameters and a body to evaluate over
them. The result is a \mcode{CompiledFunction}, an opaque object you call exactly
like any function.

\begin{codepairs}
\pair{08-compilation/basics/1}{\B{Compile} holds its arguments and returns a
\texttt{CompiledFunction}---the body is not evaluated symbolically, it is
\emph{compiled}. It echoes showing its parameters and the source body.}
\pair{08-compilation/basics/2}{Bind it to a name and it behaves like a function.
Here the parameter is declared \mcode{\{x, \_Real\}}: a machine real.}
\pair{08-compilation/basics/3}{Called on a real, it runs the bytecode and returns
a machine real.}
\pair{08-compilation/basics/4}{An exact integer argument is accepted and coerced
to the declared type---\(3\) enters as \(3.0\), so the result is \(10.0\).}
\pair{08-compilation/basics/5}{Parameters can be typed \mcode{\_Integer} \dots}
\pair{08-compilation/basics/6}{\dots or \mcode{\_Complex}; the compiler emits
integer or complex opcodes accordingly.}
\pair{08-compilation/basics/7}{And when an argument is \emph{not} a machine
number, the compiled function quietly hands the work back to the evaluator: a
symbolic \(y\) gives the ordinary symbolic answer. Compilation is an
optimisation, never a restriction on what you may pass.}
\end{codepairs}

The parameter list is the compiler's type declaration. A bare symbol means a
machine real; \mcode{\{x, \_Real\}}, \mcode{\{n, \_Integer\}}, and
\mcode{\{z, \_Complex\}} name the three scalar types explicitly; and a third
entry gives an \emph{array} rank, so \mcode{\{v, \_Real, 1\}} is a rank-1 vector
of reals and \mcode{\{m, \_Real, 2\}} a matrix (\S\ref{sec:compile-arrays}). Every
downstream instruction is chosen from these declared types, which is why the
compiler must know them up front.

\usagebox{Compile}

That last row of the example is the single most important thing to understand
about \mcode{Compile}: its \emph{fallback contract}\index{Compile!interpreter
fallback}. A \texttt{CompiledFunction} always agrees with the evaluator. If you
hand it an argument it cannot represent as a machine number, or if the body uses
something the compiler cannot lower (\S\ref{sec:compile-cliff}), it does not fail
and it does not lie---it re-runs the expression through the interpreter and
returns that. You never trade correctness for speed; at worst you fail to gain
the speed.

\section{What the compiler produces}
\label{sec:compile-bytecode}

The best way to see what compilation \emph{is} is to look at what it produces.
\B{CompilePrint} disassembles a compiled function and writes its bytecode to the
terminal:

\compileprint{08-compilation/scalar}

Read it from the top. The \emph{signature} records the parameter types and the
source body. Then three \emph{registers}: \texttt{R0} holds the argument, the
result comes back in \texttt{R1}, and one scratch register was enough. Then eight
instructions. Each is \emph{monomorphic}\index{compilation!monomorphic opcodes}:
\texttt{ADD\_RK} is ``add a real constant \(K\) to a real register,''
\texttt{MUL\_R} is ``multiply two real registers.'' There is no run-time test of
what type the operands are, because the answer was fixed when the program was
built---that erasure of per-node type dispatch is the whole source of the
speed-up. Notice too that Horner's nested form fell straight out of the source:
the compiler folded each numeric coefficient into an \emph{immediate} operand
(the \(K\) in \texttt{ADD\_RK}) rather than loading it from a register. The
program ends, as every compiled program does, with \texttt{RET}.

Not every body is straight-line. Control flow compiles to
jumps\index{compilation!control flow}: a conditional becomes a compare into a
Boolean register and a \emph{jump-if-zero} (\texttt{JZ}) around the branch not
taken, with the jump targets marked \texttt{>}.

\compileprint{08-compilation/branch}

And a counted loop compiles to a single back-edge---here \texttt{LOOP}, over
\emph{integer} opcodes (\texttt{POWI\_I}, \texttt{ADD\_I}), because every value in
this body is an integer and the ``all-Real fast path'' note is absent:

\compileprint{08-compilation/loop}

\calloutnote{Under the hood}{The register machine is a threaded-code interpreter:
on GCC and Clang each opcode's handler ends by computing the address of the next
handler and jumping straight to it (a computed \texttt{goto}), so there is no
per-instruction \texttt{switch}. Registers are a flat array of 16-byte slots on
the C stack, banked into scalar, array and ``tile'' regions, so a scalar program
allocates nothing at all per call. Elementary and special functions
(\texttt{Sin}, \texttt{Gamma}, \texttt{Erf}, \dots) do not each get their own
opcode; they share a \texttt{KERN} instruction that carries a pointer to the same
machine kernel the packed-array code uses (\S\ref{sec:ds-packed}), so a new
special function becomes compilable for free.}

You do not usually need the disassembly. What you do need, often, is a quick
answer to ``did this body actually compile, and to what?''---and for that there
is \B{CompileDiagnostics}, which type-checks a body without building the object
and reports what it found:

\begin{codepairs}
\pair{08-compilation/diagnostics/1}{The result type is \texttt{Real}; the
optimised program is three instructions, down from five before the optimiser
(constant folding, common-subexpression elimination and dead-code removal run on
the bytecode).}
\pair{08-compilation/diagnostics/2}{A trigonometric identity the compiler does
\emph{not} know: it dutifully computes both \(\sin^2\) and \(\cos^2\) in six
instructions rather than returning \(1\). The compiler is a lowering, not a
simplifier---simplify first, then compile.}
\pair{08-compilation/diagnostics/3}{Array reductions compile too:
\(\mathtt{Total}[v] + \mathtt{Max}[v]\) over a real vector is twelve
instructions, all machine code.}
\end{codepairs}

\usagebox{CompileDiagnostics}
\usagebox{CompilePrint}

\section{How fast? Compiled against the tree}
\label{sec:compile-speed}

Numbers make the case. Here is one polynomial body run over a million points two
ways: once as a \texttt{CompiledFunction}, and once as an ordinary definition
\mcode{treePoly[x\_] := \dots} that the evaluator must rewrite node by node. We
time each with \B{AbsoluteTiming}, which reports elapsed wall-clock seconds, and
take the ratio:

\mtranscript{08-compilation/speed-scalar}

\noindent The compiled loop is more than ten times faster here, for identical
output. (The exact ratio is re-measured every time this book is built and depends
on the machine; it is the \emph{order} of the thing that is stable.)

\calloutnote{Pitfall}{Measure with \B{AbsoluteTiming} (wall-clock), not \B{Timing}
(CPU seconds \emph{summed over threads}): a compiled array kernel that uses
several cores can report a \mcode{Timing} larger than the wall-clock time it
actually took, which reads backwards. And time distinct expressions, or discard a
warm-up run: the first evaluation of anything pays one-off costs (compilation
itself, cold caches) that a fair comparison should not count.}

How does that hold up against a mature numerical stack? The table below pits
Mathilda's compiled kernels against Python---vectorised NumPy where a vectorised
form exists, a plain CPython loop where one does not\calloutnote{Performance}{From
\texttt{docs/design/performance.md}, measured 2026-07-30 on an Intel Core
i9-9880H with NumPy 2.4.4 / SciPy 1.17.1 on CPython 3.11 (both linked against
Apple Accelerate). Lennard-Jones, Mandelbrot and Monte~Carlo are vectorised NumPy
and are genuine library comparisons; the logistic map has no vectorised form
(numba was not installed), so its row is ``what a Python user without a JIT
measures,'' not a library result.}:

\begin{center}
\begin{tabular}{lrr}
\hline
Benchmark & Mathilda & Python / NumPy \\
\hline
Lennard-Jones energy, 1452 bodies (all pairs) & 133.19 ms & 91.25 ms \\
Mandelbrot, \(800\times800\), 100 iterations & 769.23 ms & 980.20 ms \\
Monte Carlo \(\pi\), \(10^{7}\) samples (vectorised) & 162.36 ms & 204.53 ms \\
Logistic map, \(10^{7}\) iterations & 183.76 ms & 460.57 ms \\
\hline
\end{tabular}
\end{center}

\noindent Mathilda's compiled code wins three of the four rows and is close on the
fourth. It is worth being precise about that fourth: Lennard-Jones is where
NumPy's vectorised all-pairs form, running a hand-tuned C loop under one Python
call, comes out ahead---91~ms against 133. That is the honest shape of it. Where
the work is a genuine scalar recurrence with no array form, such as the logistic
map, the gap the other way is largest, because that is exactly the case a
vectorised library cannot help with and a bytecode VM can.

\section{The compilable subset is a cliff}
\label{sec:compile-cliff}

The compiler does not handle every function---and the way it fails is the single
most important operational fact in this chapter. The compilable
subset\index{compilation!compilable subset} is a \emph{cliff, not a slope}: if
\emph{one} head anywhere in the body cannot be lowered, the \emph{whole} body
falls back to the interpreter. Not the offending sub-expression---all of it. And
because the fallback is silent and gives the right answer, nothing tells you it
happened except the clock.

\B{CompileDiagnostics} is how you see the edge of the cliff. Give it the argument
types and a body, and it tells you whether the body compiles and, if not, the
exact innermost sub-expression that stopped it:

\begin{codepairs}
\pair{08-compilation/cliff/1}{\(\mathtt{Sin}[x] + \mathtt{Gamma}[x]\) compiles:
both heads have machine kernels.}
\pair{08-compilation/cliff/2}{Replace \texttt{Gamma} with \B{BarnesG}, which has
no machine kernel, and the whole body is rejected---the report names
\texttt{BarnesG[x]} as the culprit even though \(\sin\) beside it was perfectly
compilable.}
\pair{08-compilation/cliff/3}{The same over arrays: \B{Total} lowers, but
\B{Riffle} does not, so \(\mathtt{Total}[\mathtt{Riffle}[v, 0.]]\) does not
compile at all.}
\pair{08-compilation/cliff/4}{And an integer body: \B{MoebiusMu} has no machine
lowering, so it sinks the sum with \B{EulerPhi} regardless of the latter.}
\end{codepairs}

The lesson is a workflow. When a numeric loop is slower than you expected, do not
guess---ask \mcode{CompileDiagnostics} whether the body even compiled, fix or
remove the one head it names, and ask again. The reward for staying on the subset
is total; the penalty for stepping off it, by one head, is also total.

\calloutnote{Under the hood}{Which heads are on the subset is not a fixed list but
a moving frontier, and it is guarded. A build-time audit,
\texttt{make check-compile-coverage}, cross-checks every head that has a fast
numeric path in the interpreter against whether it also compiles, and fails the
build if a newly-added numeric builtin forgets its compiled lowering---so the
cliff can only ever recede. When an auto-compiled builtin
(\S\ref{sec:compile-auto}) falls back, setting the environment variable
\texttt{MATHILDA\_COMPILE\_DIAG=1} makes it name the head that forced it, on
standard error.}

\section{Auto-compilation}
\label{sec:compile-auto}

Most of the time the compiler works for you without being asked. When you write a
\B{Table} over a real range, or a \B{Plot}, or an \B{NIntegrate}, or a \B{Nest}
with a machine-number seed, Mathilda tries to compile the body once and run the
compiled form at every sample point or iteration. This is \emph{auto-compilation},
and it is why numeric work in Mathilda is fast without a single call to
\mcode{Compile} in sight.

The one switch that governs it is \B{\$AutoCompilation}\index{auto-compilation}.
It is \texttt{True} by default; set it to \texttt{False} and every such body goes
through the interpreter instead. Because a compiled body is contracted to give the
interpreter's answer, the switch changes speed and \emph{nothing else}---which is
exactly what makes it useful for measuring the effect. The session below turns it
off and on around the same \B{Table}, confirms the result is byte-for-byte the
same either way, and then times a million-step logistic-map iteration each way:

\mtranscript{08-compilation/autocompile}

\noindent The two tables (\texttt{r1} and \texttt{r2}) are equal: the switch is a
pure performance control, never a semantic one. The final ratio---dozens to
one---is the interpreter tax that auto-compilation removes from the
\(3.5\,x(1-x)\) recurrence. (As before, the exact figure is re-measured on each
build; its order is the point.)

Two mechanisms hide behind the one switch. An \emph{adapter} compiles a held body
once and reuses it across many sample points---this is what \B{Table}, \B{Plot},
\B{Plot3D}, \B{NIntegrate}, \B{NSum}, \B{Sum}, \B{Product} and \B{FindRoot} use.
A separate \emph{loop compiler} handles the bodies of \B{Do}, \B{While},
\B{Nest}, \B{Fold}, \B{Map} and their kin. Both are gated by the same
\mcode{\$AutoCompilation} flag, and both are conservative in the same way: an
exact-integer \mcode{Table} iterator such as \mcode{\{i, 1, n\}} is left to the
interpreter, because exactness must be preserved, while an inexact one like
\mcode{\{i, 1., n\}} is fair game.

\calloutnote{Pitfall}{Auto-compilation never wraps a machine integer
(\S\ref{sec:compile-precision}): you did not ask it to compile anything, so it may
not turn a transparent speed-up into a silently different answer. If an
auto-compiled integer computation would overflow, it abandons the compiled run and
lets the interpreter promote to a bignum---the safe default is not negotiable
behind your back.}

\section{Arrays, packing, and fusion}
\label{sec:compile-arrays}

A parameter with a rank compiles to code that operates on a whole array in one
call. Declare \mcode{\{u, \_Real, 1\}} and the body is compiled to run over the
backing buffer of a packed array (\S\ref{sec:ds-packed}), start to finish, with no
per-element expression ever built:

\mtranscript{08-compilation/arrays}

\noindent The result of a million-element call is itself a packed array
(\B{NDArrayQ} confirms it), ready to feed the next operation without unpacking.

What makes the array path fast is worth stating precisely, because it is not what
one first assumes. Look again at the disassembly of exactly this body:

\compileprint{08-compilation/array}

\noindent The element-wise chain \(u^2 + 1\) was \emph{fused}\index{compilation!fusion}
into a single strip-mined pass: one parallel loop (\texttt{APAR}) walks the buffer
in 64-wide tiles, and within each tile the square and the add happen on tile
registers (\texttt{VPOWI\_R}, \texttt{VADD\_R}) before the result is written back
(\texttt{VSTORE\_R}). There is no intermediate array holding \(u^2\). Array speed
in Mathilda does not come from ``removing interpretation''---the packed
interpreter is already fast, because it too calls machine kernels
(\S\ref{sec:ds-packed}). It comes from \emph{removing per-element dispatch and
intermediate buffers}, which is why compiling a length-16 vector barely helps
while compiling a long one, or a longer chain of operations, helps a great deal.

\calloutnote{Performance}{This is the same lesson the packed-array chapter
teaches from the other side. Packing decides whether the data is a dense buffer at
all; compilation, and the fusion it enables, decides how few passes over that
buffer the arithmetic takes. The two compose: a compiled body over packed
arguments is the fast path, and it is the one \B{Table}, \B{Plot} and the numeric
solvers of Chapter~\ref{ch:datastructures}'s sequel put you on automatically.}

\section{Machine integers and arbitrary precision}
\label{sec:compile-precision}

By default the compiler's integers are machine \texttt{int64}s and its reals are
machine \texttt{double}s---that is the point, and it is what makes the code fast.
But two questions follow immediately: what happens when a machine integer
overflows, and what if machine precision is not enough? Mathilda has a careful
answer to each.

\begin{codepairs}
\pair{08-compilation/precision/1}{\(3000000^3\) is far past the \texttt{int64}
range. The compiled code detects the overflow, abandons the call, and the
interpreter re-runs it and promotes to a bignum \dots}
\pair{08-compilation/precision/2}{\dots giving exactly the interpreter's answer.
Overflow is caught, never wrapped: the compiled result is always the true one.}
\pair{08-compilation/precision/3}{For more digits than a \texttt{double} holds,
\mcode{WorkingPrecision -> n} compiles the arithmetic in extended precision (MPFR)
at \(n\) digits---here \(\sqrt{2}\) to forty \dots}
\pair{08-compilation/precision/4}{\dots identical to what the interpreter computes
at the same precision.}
\pair{08-compilation/precision/5}{The diagnostics confirm the result type is now
\texttt{MPFRReal}, not \texttt{Real}.}
\pair{08-compilation/precision/6}{Alternatively \mcode{"BigIntegers" -> True}
keeps integer arithmetic exact (GMP) \emph{inside} the compiled code, so the cube
stays compiled and correct rather than bailing out on overflow \dots}
\pair{08-compilation/precision/7}{\dots with result type \texttt{BigInteger}.}
\end{codepairs}

There is a third mode, for the rare case where you know your integers stay within
range and want to skip even the overflow check: \mcode{RuntimeOptions -> "Speed"}.
It removes the per-operation check, so an overflowing result is left wrapped
modulo \(2^{64}\) instead of being promoted:

\begin{codepairs}
\pair{08-compilation/precision/8}{The same cube, now in \texttt{"Speed"} mode: no
promotion, so the answer is the low 64 bits of the true result---a different, and
here plainly wrong, number.}
\pair{08-compilation/precision/9}{Which is exactly \(3000000^3 \bmod 2^{64}\): the
machine word has simply wrapped.}
\end{codepairs}

\noindent This is the one place in the chapter where the compiled answer is
\emph{not} contracted to the true one, and it is off by default for that reason.
Reach for it only when you have measured that the check matters and proven that
your values cannot overflow.

\calloutnote{Under the hood}{The extended-precision modes are deliberately narrow:
\mcode{WorkingPrecision} and \mcode{"BigIntegers"} cover straight-line arithmetic
and the elementary functions on scalars---enough for the sample-point bodies the
numeric solvers of \S\ref{sec:numcalc} compile at high precision---but not arrays
or control flow, which stay on the machine path. The machine path itself is
untouched by these options: switch them off and the bytecode is identical to the
last byte.}

\section{When to compile}
\label{sec:compile-when}

The rule of thumb is short. If a numeric expression is evaluated many times---in a
loop, over an array, at every abscissa of an integration---compilation is a large,
free win, and most of the time (\S\ref{sec:compile-auto}) you get it without
asking. Reach for \mcode{Compile} explicitly when you want to build the fast
function once and reuse it, when you want to hand it to \B{Map} or a solver, or
when you simply want to see with \mcode{CompileDiagnostics} and \mcode{CompilePrint}
what the machine is doing. Do \emph{not} compile a symbolic computation, a
one-shot evaluation, or a body full of heads off the subset---there the fallback
gives you nothing but the overhead of trying.

The threads pulled together here run through the rest of the book. The packed
arrays the compiler operates on are Chapter~\ref{ch:datastructures}; the numeric
solvers that auto-compile their integrands and right-hand sides are
\S\ref{sec:numcalc}; and the evaluator whose generality the compiler trades away
for speed is the subject of the internals to come.
